% =============================================================================
%  Part I -- Overview: the assumptions, and the ideas of each scheme
% =============================================================================
\landscapeopen
\section*{Part I --- Overview}

% ---------------------------------------------------------------------------
%  Table I.0  glyph matrix
% ---------------------------------------------------------------------------
\subsection*{Table I.0 --- At a glance: which assumption does each scheme drop?}
\begin{center}
\small
\setlength{\tabcolsep}{6pt}
\renewcommand{\arraystretch}{1.15}
\newcommand{\rl}[2]{\rid{#1}\ #2}
\begin{tabular}{@{}l
  >{\columncolor{tMYNN}}c c c c @{\hspace{7pt}}|@{\hspace{7pt}}
  >{\columncolor{tGXVIII}}c >{\columncolor{tGXIX}}c c >{\columncolor{tAPPROX}}c >{\columncolor{tFULL}}c c c@{}}
 & \multicolumn{4}{c}{\footnotesize\textit{1D family}} & \multicolumn{7}{c}{\footnotesize\textit{3D family}} \\[-0.4ex]
 & \hdr{mynnrun} & \hdr{ysu} & \hdr{shh} & \hdr{boulac} & \hdr{gxviii} & \hdr{gxix} & \hdr{ito} & \hdr{approx} & \hdr{full} & \hdr{sms} & \hdr{les} \\
\midrule
\rowshade{A1}\rl{A1}{all energetic eddies are subgrid} & \G{A1}{mynnrun} & \G{A1}{ysu} & \G{A1}{shh} & \G{A1}{boulac} & \G{A1}{gxviii} & \G{A1}{gxix} & \G{A1}{ito} & \G{A1}{approx} & \G{A1}{full} & \G{A1}{sms} & \G{A1}{les} \\
\rowshade{A2}\rl{A2}{cell mean = ensemble mean} & \G{A2}{mynnrun} & \G{A2}{ysu} & \G{A2}{shh} & \G{A2}{boulac} & \G{A2}{gxviii} & \G{A2}{gxix} & \G{A2}{ito} & \G{A2}{approx} & \G{A2}{full} & \G{A2}{sms} & \G{A2}{les} \\
\addlinespace[2pt]
\rowshade{B1}\rl{B1}{no horizontal shear production} & \G{B1}{mynnrun} & \G{B1}{ysu} & \G{B1}{shh} & \G{B1}{boulac} & \G{B1}{gxviii} & \G{B1}{gxix} & \G{B1}{ito} & \G{B1}{approx} & \G{B1}{full} & \G{B1}{sms} & \G{B1}{les} \\
\rowshade{B2}\rl{B2}{no gradients of $W$} & \G{B2}{mynnrun} & \G{B2}{ysu} & \G{B2}{shh} & \G{B2}{boulac} & \G{B2}{gxviii} & \G{B2}{gxix} & \G{B2}{ito} & \G{B2}{approx} & \G{B2}{full} & \G{B2}{sms} & \G{B2}{les} \\
\rowshade{B3}\rl{B3}{no horizontal flux divergence} & \G{B3}{mynnrun} & \G{B3}{ysu} & \G{B3}{shh} & \G{B3}{boulac} & \G{B3}{gxviii} & \G{B3}{gxix} & \G{B3}{ito} & \G{B3}{approx} & \G{B3}{full} & \G{B3}{sms} & \G{B3}{les} \\
\rowshade{B4}\rl{B4}{TKE stays in its column} & \G{B4}{mynnrun} & \G{B4}{ysu} & \G{B4}{shh} & \G{B4}{boulac} & \G{B4}{gxviii} & \G{B4}{gxix} & \G{B4}{ito} & \G{B4}{approx} & \G{B4}{full} & \G{B4}{sms} & \G{B4}{les} \\
\rowshade{B5}\rl{B5}{vertical and horizontal mixing separable} & \G{B5}{mynnrun} & \G{B5}{ysu} & \G{B5}{shh} & \G{B5}{boulac} & \G{B5}{gxviii} & \G{B5}{gxix} & \G{B5}{ito} & \G{B5}{approx} & \G{B5}{full} & \G{B5}{sms} & \G{B5}{les} \\
\rowshade{B6}\rl{B6}{flat, homogeneous surface (MOST)} & \G{B6}{mynnrun} & \G{B6}{ysu} & \G{B6}{shh} & \G{B6}{boulac} & \G{B6}{gxviii} & \G{B6}{gxix} & \G{B6}{ito} & \G{B6}{approx} & \G{B6}{full} & \G{B6}{sms} & \G{B6}{les} \\
\rowshade{B7}\rl{B7}{model surfaces are horizontal} & \G{B7}{mynnrun} & \G{B7}{ysu} & \G{B7}{shh} & \G{B7}{boulac} & \G{B7}{gxviii} & \G{B7}{gxix} & \G{B7}{ito} & \G{B7}{approx} & \G{B7}{full} & \G{B7}{sms} & \G{B7}{les} \\
\addlinespace[2pt]
\rowshade{C1}\rl{C1}{fluxes local and down-gradient} & \G{C1}{mynnrun} & \G{C1}{ysu} & \G{C1}{shh} & \G{C1}{boulac} & \G{C1}{gxviii} & \G{C1}{gxix} & \G{C1}{ito} & \G{C1}{approx} & \G{C1}{full} & \G{C1}{sms} & \G{C1}{les} \\
\rowshade{C2}\rl{C2}{one $K$ for all directions} & \G{C2}{mynnrun} & \G{C2}{ysu} & \G{C2}{shh} & \G{C2}{boulac} & \G{C2}{gxviii} & \G{C2}{gxix} & \G{C2}{ito} & \G{C2}{approx} & \G{C2}{full} & \G{C2}{sms} & \G{C2}{les} \\
\rowshade{C3}\rl{C3}{flux anti-parallel to the gradient} & \G{C3}{mynnrun} & \G{C3}{ysu} & \G{C3}{shh} & \G{C3}{boulac} & \G{C3}{gxviii} & \G{C3}{gxix} & \G{C3}{ito} & \G{C3}{approx} & \G{C3}{full} & \G{C3}{sms} & \G{C3}{les} \\
\rowshade{C4}\rl{C4}{eddy size from the vertical structure} & \G{C4}{mynnrun} & \G{C4}{ysu} & \G{C4}{shh} & \G{C4}{boulac} & \G{C4}{gxviii} & \G{C4}{gxix} & \G{C4}{ito} & \G{C4}{approx} & \G{C4}{full} & \G{C4}{sms} & \G{C4}{les} \\
\rowshade{C5}\rl{C5}{one scale for transport and dissipation} & \G{C5}{mynnrun} & \G{C5}{ysu} & \G{C5}{shh} & \G{C5}{boulac} & \G{C5}{gxviii} & \G{C5}{gxix} & \G{C5}{ito} & \G{C5}{approx} & \G{C5}{full} & \G{C5}{sms} & \G{C5}{les} \\
\rowshade{C6}\rl{C6}{vertical Ri only; turbulence dies at $\Rfc$} & \G{C6}{mynnrun} & \G{C6}{ysu} & \G{C6}{shh} & \G{C6}{boulac} & \G{C6}{gxviii} & \G{C6}{gxix} & \G{C6}{ito} & \G{C6}{approx} & \G{C6}{full} & \G{C6}{sms} & \G{C6}{les} \\
\rowshade{C7}\rl{C7}{local equilibrium of all moments but $\qsq$} & \G{C7}{mynnrun} & \G{C7}{ysu} & \G{C7}{shh} & \G{C7}{boulac} & \G{C7}{gxviii} & \G{C7}{gxix} & \G{C7}{ito} & \G{C7}{approx} & \G{C7}{full} & \G{C7}{sms} & \G{C7}{les} \\
\rowshade{C8}\rl{C8}{buoyancy acts only through $\avg{w\theta_v}$} & \G{C8}{mynnrun} & \G{C8}{ysu} & \G{C8}{shh} & \G{C8}{boulac} & \G{C8}{gxviii} & \G{C8}{gxix} & \G{C8}{ito} & \G{C8}{approx} & \G{C8}{full} & \G{C8}{sms} & \G{C8}{les} \\
\rowshade{C9}\rl{C9}{universal closure constants} & \G{C9}{mynnrun} & \G{C9}{ysu} & \G{C9}{shh} & \G{C9}{boulac} & \G{C9}{gxviii} & \G{C9}{gxix} & \G{C9}{ito} & \G{C9}{approx} & \G{C9}{full} & \G{C9}{sms} & \G{C9}{les} \\
\rowshade{C10}\rl{C10}{pressure terms: transport absorbed, Rotta} & \G{C10}{mynnrun} & \G{C10}{ysu} & \G{C10}{shh} & \G{C10}{boulac} & \G{C10}{gxviii} & \G{C10}{gxix} & \G{C10}{ito} & \G{C10}{approx} & \G{C10}{full} & \G{C10}{sms} & \G{C10}{les} \\
\bottomrule
\end{tabular}
\end{center}
{\footnotesize Reading the matrix. The shaded rows are what no scheme touches: the ensemble
average, the flat surface layer, local equilibrium, universal constants and the pressure
terms --- the frontier left open by all of them. Down the 3D family the number of open circles
grows from the Goger ports (one source term) to the full closure (production, divergences,
transport and the tensor itself); SMS-3DTKE and the LES limit drop the same geometric
assumptions but keep the eddy-viscosity form. The lone open circle in the MYNN-as-run column is
the removed critical Richardson number --- the control run has no stability cutoff, while the
approximate 3D closure still carries one.\par}
\landscapeclose

% ---------------------------------------------------------------------------
%  Table I.1  the assumptions
% ---------------------------------------------------------------------------
\clearpage
\subsection*{Table I.1 --- The assumptions, where MYNN 2.5 uses them, and why they fail}
\footnotesize
\setlength{\tabcolsep}{4pt}
\renewcommand{\arraystretch}{1.12}
\begin{xltabular}{\linewidth}{@{}>{\hsize=1.05\hsize\raggedright\arraybackslash}X
  >{\hsize=0.85\hsize\raggedright\arraybackslash}X
  >{\hsize=1.15\hsize\raggedright\arraybackslash}X
  >{\hsize=0.95\hsize\raggedright\arraybackslash}X@{}}
\toprule
\textbf{Assumption} (smallest statement) & \textbf{Where MYNN 2.5 uses it} & \textbf{Why it fails in mountains} & \textbf{Why it fails in the grey zone} \\
\midrule
\endfirsthead
\toprule
\textbf{Assumption} (continued) & \textbf{Where MYNN 2.5 uses it} & \textbf{Why it fails in mountains} & \textbf{Why it fails in the grey zone} \\
\midrule
\endhead
\bottomrule
\endlastfoot
\multicolumn{4}{@{}l}{\textit{Family A --- what the grid cell is}} \\
\addlinespace[2pt]
\rid{A1} Every energy-containing eddy is smaller than the grid cell; the closure carries the whole spectrum and the grid spacing appears nowhere in it. \newline $\Delta \gg l,\ z_i$; \ $\partial(\text{closure})/\partial\Delta = 0$ &
$L$, $S_M$, $S_H$ and $\varepsilon$ contain no $\Delta$; the resolved flow is taken as non-turbulent, so the grid-mean shear enters $P_s$ as if it were the ensemble-mean shear. &
At $\dx = 500$~m the daytime valley mixed layer ($z_i \approx 1$--$2$~km) is partly resolved: the up-valley jet and the slope-flow branches are resolved circulations whose shear the scheme also parameterises --- double counting. In this project's convective morning 85--91\,\% of the kinetic energy at 43~m was resolved. \rg{day}\ev{run} &
For $\Delta \approx z_i$ the same eddies appear as resolved motion and as parameterised flux; the parameterised share should follow the partition of \textcite{honnert2011} --- a function of $\Delta/(z_i + h_c)$ alone, one half at 0.2, i.e.\ 200~m for a 1~km mixed layer --- while a non-scale-aware scheme keeps it at 100\,\% and weakens or delays the resolved convection \parencite{ching2014}. What the grid does resolve sits at $6$--$8\Delta$, the effective resolution \parencite{skamarock2004,zhou2014}: at 500~m the smallest trustworthy eddy is 3--4~km wide. At night the dominant eddy shrinks to $\sim$100~m and a 500~m grid is mesoscale again \parencite{zhou2014}: the stable regime is not a grey-zone problem. Every partition result is flat-terrain free convection; no reference partition exists for the stable regime. \rg{day}\ev{inf} \\
\addlinespace[3pt]
\rid{A2} The grid-cell average is an ensemble average: fluctuations average to zero, the mean of a product is the product of the means plus the covariance, and the closure constants are ensemble constants. \newline $\avg{u'} = 0$, $\avg{U u'} = 0$ &
The Reynolds-averaged budget of $\qsq$; the Mellor--Yamada constants were fitted to ensemble statistics of neutral and convective flat-terrain flows. &
A 500~m cell on a slope holds the slope-wind layer and the free valley air side by side: its mean is a mixture of two flows, not the ensemble of one. \rg{both}\ev{inf} &
The average is a spatial filter over a few eddies \parencite{wyngaard2004}: Leonard and cross terms appear, statistics fluctuate from cell to cell, and constants calibrated on ensembles need not hold. \rg{day}\ev{inf} \\
\addlinespace[4pt]
\multicolumn{4}{@{}l}{\textit{Family B --- where the gradients and the fluxes are}} \\
\addlinespace[2pt]
\rid{B1} Horizontal gradients of the mean wind and of scalars produce no turbulence. \newline $P_s = -\avg{uw}\,\pd{U}{z} - \avg{vw}\,\pd{V}{z}$ (two of nine terms); $\pd{\Theta}{x} = \pd{\Theta}{y} = 0$ in the fluxes &
$P_s = L q S_M\,[(\pd{U}{z})^2 + (\pd{V}{z})^2]$; the stability functions see only the vertical shear through $G_M = (L^2/q^2)[(\pd{U}{z})^2 + (\pd{V}{z})^2]$. &
The up-valley jet has a cross-valley shear $\pd{u}{y}$ comparable to $\pd{u}{z}$: horizontal shear reaches 20--30\,\% of the bulk shear at 500 and 100~m \parencite{freddi2026}, and horizontal shear production is observable at the valley floor in the afternoon \parencite{goger2018}; at night drainage and katabatic flows add horizontal shear that the vertical Richardson number does not see, and in weak-wind drainage flow the directional shear, not the speed shear, sets $u_*$ \parencite{mahrt2017}. Already without an along-valley jet, the idealised-valley LES of \textcite{schmidli2013} has a horizontal turbulent heat flux as large as the vertical one and shear production from the circulation-cell gradients over slopes and ridge. \rg{both}\ev{obs} &
Resolved convective structures have $\pd{U}{x} \sim \pd{U}{z}$; the six missing terms are of the order of the two retained ones --- the eddies are "neither fully parameterised nor resolved" \parencite{juliano2022}. \rg{day}\ev{inf} \\
\addlinespace[3pt]
\rid{B2} The mean vertical velocity has no gradients that matter. \newline $\pd{W}{x_j} = 0$ in the closure &
$W$ is absent from $G_M$, $G_H$ and from the fluxes; only $\avg{uw}$, $\avg{vw}$, $\avg{w\theta}$ exist. &
On a 30$^\circ$ slope the slope-parallel flow has a vertical component comparable to its horizontal one: $\pd{W}{x}$ and $\pd{W}{z}$ are of the order of the wind shear, and $-\avg{uw}\,\pd{W}{x}$ pairs with the retained $-\avg{uw}\,\pd{U}{z}$. \rg{both}\ev{inf} &
Resolved updraughts of the convective boundary layer have $\pd{W}{z} \sim w_*/z_i$, the same order as the shear the scheme keeps. \rg{day}\ev{inf} \\
\addlinespace[3pt]
\rid{B3} Turbulent fluxes have no horizontal divergence: the mean flow feels turbulence only through $\pd{\avg{\cdot\,w}}{z}$. \newline $\pd{U}{t}|_{turb} = -\pd{\avg{uw}}{z}$, $\pd{\Theta}{t}|_{turb} = -\pd{\avg{w\theta}}{z}$ &
Implicit vertical diffusion with $K_M = L q S_M$, $K_H = L q S_H$; horizontal mixing is left to WRF's numerical diffusion. &
Fluxes differ between the slope-wind layer and the valley core within one grid cell; the MAP-Riviera TKE budgets closed only with transport terms \parencite{rotach2004,weigel2004}. \rg{both}\ev{obs} &
$\pd{\avg{uu}}{x} \sim \pd{\avg{uw}}{z}$ when $\Delta \sim z_i$; the third of Juliano's three ingredients, without which the spurious secondary circulations survive. \rg{day}\ev{inf} \\
\addlinespace[3pt]
\rid{B4} TKE is a property of the column: neither advected by the resolved wind nor diffused horizontally. \newline $D\qsq/Dt \to \pd{\qsq}{t}$; transport $\pd{}{z}(L q S_q\,\pd{\qsq}{z})$ only &
The $\qsq$ equation carries vertical transport only; $\qsq$ is not a transported model variable. &
The valley wind (5--10~m\,s$^{-1}$) carries TKE produced at the slope--jet interface across a 500~m cell in one to two minutes, faster than the eddy turnover $L/q$; the Riviera budgets needed advection \parencite{rotach2004}; near the surface \textcite{goger2018} found it minor. \rg{both}\ev{obs} &
Resolved eddies advect subgrid TKE, and horizontal diffusion of TKE matters once $\Delta \sim L$. \rg{day}\ev{inf} \\
\addlinespace[3pt]
\rid{B5} Vertical and horizontal mixing are separable: the scheme mixes vertically from its TKE, a numerical Smagorinsky operator mixes horizontally with the grid as length scale, and neither exchanges energy with the other. \newline $K_v = L q S_M$; $K_h = (c_s \Delta_h)^2 |D_h|$; $P_s$ contains no $K_h$ term &
In WRF the scheme runs with \texttt{diff\_opt=2}, \texttt{km\_opt=4}: $K_h = (0.25\,\Delta_h)^2 |D_h|$ acts on the horizontal deformation only, knows no stability and returns nothing to $\qsq$. &
The numerical operator stands exactly where the physical horizontal flux divergence belongs \parencite{juliano2022}, and the energy it removes from the resolved slope and valley flows vanishes instead of feeding turbulence. The operator was calibrated at $\Delta = 555$~km, and its author already suspected it of playing "too large a role" \parencite{smagorinsky1963}. \rg{both}\ev{inf} &
$K_h$ and $K_v$ are physically of the same order at $\Delta \sim z_i$; treating them separately is the "energetic inconsistency" that motivates SMS-3DTKE \parencite{zhang2018}. Worse, vertical-only diffusion removes the short-wave cutoff of the Rayleigh instability of the scheme's own superadiabatic layer: spurious cells at $\lambda \approx 8\Delta$ grow with an e-folding time of 20~min at 500~m against 3~h at 5~km, most for local (down-gradient) schemes \parencite{ching2014}. \rg{day}\ev{inf} \\
\addlinespace[3pt]
\rid{B6} The lower boundary is a flat, homogeneous, stationary surface: Monin--Obukhov similarity holds, fluxes are normal to a horizontal surface, one roughness length per cell. \newline $u_*$, $\theta_*$ from $\phi_m(z/L_{\mathrm{MO}})$, $\phi_h(z/L_{\mathrm{MO}})$ &
Lower boundary from the surface-layer scheme; the wall value $\qsq = B_1^{2/3} u_*^2$ is the level-2 equilibrium of a neutral surface layer. &
On a slope there is no constant-flux layer: the katabatic jet maximum sits at 2--10~m inside it and the gradients change sign; in the Inn Valley the fluxes are constant with height for momentum in only 8--28\,\% of the daytime periods, least on the steepest slope \parencite{lehner2021}; MOST scaling holds only for some anisotropy classes \parencite{stiperski2017}, and over a hill only below a third of the inner layer, with two Obukhov lengths for the two components of gravity \parencite{kaimal1994}. The surface layer does not represent the column: at Innsbruck the surface is stable 4.7\,\% of the time while the valley air above is often stably stratified \parencite{ward2022}. The functions themselves were fitted to 34 runs over 2400~m of uniform Kansas stubble, with a stable branch that reaches its critical Richardson number at 0.21 \parencite{businger1971}. Within one 500~m cell slope and floor coexist. \rg{both}\ev{obs} &
The ensemble flux--gradient relation of MOST is applied to the instantaneous local gradients under resolved eddies --- the classical wall-model problem of LES. \rg{day}\ev{inf} \\
\addlinespace[3pt]
\rid{B7} The coordinate directions are the physical ones: a derivative along the model surface is a horizontal derivative. \newline $\pd{}{x}|_\eta = \pd{}{x}|_z + (\pd{z}{x})_\eta\,\pd{}{z}$ &
Not needed: the scheme has no horizontal derivative. The assumption becomes live the moment B1--B3 are dropped. &
On a 30$^\circ$ slope the metric term is as large as the raw derivative; without correction the coordinate manufactures strain and mixes along sloping surfaces \parencite{zangl2002}. \rg{both}\ev{inf} &
Same magnitude problem for the resolved gradients a 3D closure consumes. \rg{day}\ev{inf} \\
\addlinespace[4pt]
\multicolumn{4}{@{}l}{\textit{Family C --- how the fluxes follow from the gradients}} \\
\addlinespace[2pt]
\rid{C1} Fluxes are local and down-gradient: each flux is proportional to the local vertical gradient of its own mean quantity. \newline $\avg{w\theta} = -K_H\,\pd{\Theta}{z}$, $\avg{uw} = -K_M\,\pd{U}{z}$ &
The K-form of level 2.5; level 3 would add the countergradient term $\Gamma_\theta$. &
Turbulence produced aloft --- at the slope--valley interface, at elevated inversions --- transports heat and momentum without or against the local gradient \parencite{rotach2004,lehner2018}; at the nose of the nocturnal jet the shear vanishes while the stress does not. \rg{both}\ev{obs} &
The nonlocal (plume) part of the flux is what becomes resolved first; a local K-model has no such part to hand over \parencite{shin2015}. Once eddies are resolved the local scheme no longer needs to sustain a superadiabatic profile: at 500~m MYNN leaves over 80\,\% of the heat flux to resolved motions, heat compensating between subgrid and resolved while momentum does not \parencite{shin2016}. \rg{day}\ev{inf} \\
\addlinespace[3pt]
\rid{C2} The eddy viscosity is a scalar: one magnitude for all directions (one $K$, one $L$). &
$K_M$ serves $\avg{uw}$ and $\avg{vw}$ alike; no horizontal $K$ exists inside the scheme. &
Shear-generated valley turbulence is anisotropic --- two-component, pancake-like in stable air \parencite{stiperski2017}; the horizontal eddy scale is set by the valley width, the vertical by the layer depth: at least two length scales \parencite{goger2019}. Even over flat terrain the shear-driven layer is $\avg{u^2}$-dominated near the surface while the convective one is $\avg{w^2}$-dominated aloft, with twice the $uw$ correlation \parencite{moeng1994}. \rg{both}\ev{obs} &
With $\Delta_h = 500$~m and $\Delta z = 20$~m the filter itself is anisotropic; subgrid eddies cannot be isotropic in such a box. \rg{day}\ev{inf} \\
\addlinespace[3pt]
\rid{C3} The flux is anti-parallel to the gradient of its own quantity (Boussinesq alignment). \newline $\avg{u_i u_j} \propto -(\pd{U_i}{x_j} + \pd{U_j}{x_i})$ &
Implied by the K-form: the stress tensor is a scalar times the strain. &
Where the slope wind turns into the valley wind the principal axes of the stress tensor rotate away from the vertical gradient: the lateral momentum flux is not anti-parallel to it \parencite{zardi2013,lehner2021}, and locally negative production is admissible; the observed Inn Valley up-valley jet veers clockwise with height through its 200--400~m core \parencite{babic2021}. \rg{both}\ev{inf} &
Strain from resolved eddies rotates faster than the relaxation time $L/q$; the algebraic equilibrium behind the alignment has no basis. \rg{day}\ev{inf} \\
\addlinespace[3pt]
\rid{C4} The eddy size is a vertical-structure quantity: distance to the ground, depth of the turbulent layer, stratification; it knows neither the terrain geometry nor the grid. \newline $1/L = 1/L_S + 1/L_T + 1/L_B$ &
$L_S = \kappa z/(1 + 2.7\zeta)$ (stable) or $\kappa z (1 - 100\zeta)^{0.2}$ (unstable), $\zeta = z/L_{\mathrm{MO}}$; $L_T = 0.23\int q z\,dz/\int q\,dz$; $L_B = [1 + 5 (q_c/L_T N)^{1/2}]\,q/N$. &
"Distance to the ground" is ambiguous on a slope; $z_i$ is ill-posed in a valley with elevated inversions and multiple layers \parencite{lehner2018} --- over mountains the boundary-layer height has no universal definition, and the one fit for parametrising fluxes (the inversion base) differs from the one describing the layer the thermally driven flows reach \parencite{dewekker2015}; the bulk-Richardson height behind several schemes was fitted at 0.22 in its classical form and 0.25 with the friction term, on flat Cabauw nights \parencite{vogelezang1996}. The horizontal eddy scale is the valley width or the flow, which is why $\dx$ was a "lucky choice" \parencite{goger2018,goger2019}. At 750~m in a stable cold pool the length scale outweighed the horizontal fluxes in the 3D closure's response, and it is blamed for the full solution's instability \parencite{arthur2022}. \rg{both}\ev{inf} &
$L$ must not exceed the filter scale (LES limit $l \to \Delta$); a mesoscale $L$ of hundreds of metres on a 500~m grid double counts the eddies the grid resolves \parencite{honnert2011,ito2015}. Conversely $\Delta$ is not an eddy size: the smallest resolved eddy is $\sim 7\Delta$ \parencite{skamarock2004}, so $l_h = 0.2\Delta$ or $\Delta$ is a grid length, and "the use of grid spacing as the relevant turbulent length-scale fails at grey zone resolutions" \parencite{efstathiou2015}. \rg{day}\ev{inf} \\
\addlinespace[3pt]
\rid{C5} One length scale governs both transport and dissipation. \newline $\varepsilon = q^3/(B_1 L)$ with the same $L$ as in $K = L q S$ &
The dissipation is diagnosed from the master length scale; production and decay refer to the same eddy. &
Load-bearing rather than optional: this project's nocturnal runaway of $\qsq$ came from limiting $l$ in the flux closure but not in $\varepsilon$; one master scale removed it. \rg{night}\ev{run} &
Dissipation is nearly filter-independent while $q^3$ shrinks with $\Delta$, so $L_\varepsilon = q^3/\varepsilon$ must shrink with the grid \parencite{ito2015}; in Deardorff's model $C_\varepsilon$ depends on $l/\Delta$. \rg{day}\ev{inf} \\
\addlinespace[3pt]
\rid{C6} Stability is measured by the vertical gradient Richardson number alone, and turbulence cannot survive beyond a critical value. \newline $\Ri = N^2/[(\pd{U}{z})^2 + (\pd{V}{z})^2]$; $S_M = S_H = 0$ for $\Rf \ge \Rfc$ &
Level-2 functions $S_M(\Ri)$, $S_H(\Ri)$ with $\Rfc = \gamma_1/(\gamma_1 + \gamma_2) \approx 0.30$ for the NN09 constants, carried to level 2.5 by the Helfand--Labraga rescaling; $L_B \propto q/N$. &
The nocturnal valley atmosphere is turbulent at $\Ri \gg 1$: in this project 69\,\% of the nocturnal valley cells sit above $\Ri = 0.3$, where the 3D closure's deficit against MYNN is flat \ev{run}; intermittency, wave--turbulence interaction and drainage-flow shear \parencite{lehner2023}; horizontal shear of drainage and katabatic flows lowers the effective Richardson number, which the vertical $\Ri$ does not see. \rg{night}\ev{obs} &
Not a grey-zone problem at 500~m: the stable eddies are subgrid. The failure is terrain, not resolution. \rg{night} \\
\addlinespace[3pt]
\rid{C7} All second moments except $\qsq$ are in local equilibrium with the mean gradients: their tendencies and transport are neglected (level-2.5 truncation), and the stability functions come from the level-2 balance $P_s + P_b = \varepsilon$. &
$S_M$, $S_H$ from level 2, corrected for disequilibrium by $\alpha_c = q/q_{eq}$ where $q < q_{eq}$ \parencite{helfand1988}. &
Transitions (evening, drainage onset), intermittency and advection make turbulence non-local to its production; the equilibrium time $L/q$ (minutes) is not short against them. \rg{both}\ev{inf} &
Resolved strain changes faster than $L/q$: $|S|\,L/q \gg 1$, the regime the 3D closure's strain cap has to guard against. \rg{day}\ev{run} In the convective regime the third-moment transport of TKE is as large as production and dissipation, so no level-2 balance holds locally \parencite{moeng1994}. \rg{day}\ev{inf} \\
\addlinespace[3pt]
\rid{C8} Buoyancy acts on the TKE budget only, through the vertical heat flux; it neither redistributes energy among the stress components nor creates horizontal heat fluxes. \newline $P_b = \beta g\,\avg{w\theta_v}$; $\avg{u\theta_v} = \avg{v\theta_v} = 0$ &
$P_b = -L q S_H G_H$ with $G_H = -(L^2/q^2)\,\beta g\,\pd{\Theta_v}{z}$. &
The surface heat flux on a slope is slope-normal: its horizontal component has no home in the scheme, although it drives the thermal structure of the slope-wind layer; in stable air buoyancy damps $\avg{w^2}$ selectively --- the pancake anisotropy \parencite{stiperski2017}. \rg{both}\ev{inf} &
Resolved thermals carry the horizontal heat flux; its subgrid part remains unrepresented. \rg{day}\ev{inf} \\
\addlinespace[3pt]
\rid{C9} The closure constants are universal: one set, calibrated on flat neutral or convective boundary layers, holds everywhere. \newline $(A_1, A_2, B_1, B_2, C_1, C_2, C_3, C_5)$ &
NN09: $A_1 = 1.18$, $A_2 = 0.665$, $B_1 = 24.0$, $B_2 = 15.0$, $C_1 = 0.137$, $C_2 = 0.75$, $C_3 = 0.352$, $C_5 = 0.2$, tuned to large-eddy simulations of flat convective and stable layers. &
The constants encode return-to-isotropy and dissipation rates of flat-terrain turbulence; the anisotropy classes of complex terrain differ \parencite{stiperski2017}; \textcite{kosovic2020} conclude that "additional parameter tuning may be necessary", and the schemes compared here already run three different sets. \rg{both}\ev{inf} &
Ensemble constants versus filter constants: the LES models' $C_\varepsilon$ depends on $l/\Delta$; the same value cannot serve both limits. \rg{day}\ev{inf} \\
\addlinespace[3pt]
\rid{C10} Pressure--velocity correlations are simple: pressure transport of TKE is absorbed into the down-gradient TKE transport, and the pressure--strain term is a return to isotropy plus an isotropisation of production. \newline $\avg{u_j(e + p/\rho)} = -K_q\,\pd{e}{x_j}$; $\Pi_{ij} = -\frac{q}{3A_1 L}(\avg{u_i u_j} - \tfrac{\delta_{ij}}{3}\qsq) + C_1 \qsq\,\pd{U_i}{x_j}$ &
$S_q = 3 S_M$ in the $\qsq$ transport; Rotta's hypothesis with $A_1$, $C_1$ behind $S_M$, $S_H$. &
"Taking the pressure correlation term to be small may be a dangerous assumption" in complex terrain \parencite{goger2018}; pressure work by slope flows and gravity waves redistributes energy laterally. \rg{both}\ev{inf} &
The pressure term carries the resolved--subgrid energy exchange (the "rapid" part); it is the least known term of the budget. \rg{day}\ev{inf} \\
\end{xltabular}
\normalsize

% ---------------------------------------------------------------------------
%  Table I.2a  ideas, 1D family
% ---------------------------------------------------------------------------
\landscapeopen
\subsection*{Table I.2a --- How each scheme handles the assumption: the 1D family}
{\footnotesize The idea only; constants, equation numbers and code lines are in Part~II. \par}
\vspace{0.5ex}
\scriptsize
\setlength{\tabcolsep}{4pt}
\renewcommand{\arraystretch}{1.1}
\begin{xltabular}{\linewidth}{@{}>{\raggedright\arraybackslash}p{2.9cm}
  >{\columncolor{tMYNN}\raggedright\arraybackslash}X
  >{\raggedright\arraybackslash}X >{\raggedright\arraybackslash}X >{\raggedright\arraybackslash}X@{}}
\toprule
 & \hd{mynnrun} & \hdn{YSU} & \hdn{Shin--Hong 2015} & \hdn{BouLac} \\
\midrule
\endfirsthead
\toprule
 & \hd{mynnrun} & \hdn{YSU} & \hdn{Shin--Hong 2015} & \hdn{BouLac} \\
\midrule
\endhead
\bottomrule
\endlastfoot
\rowshade{A1}\rid{A1} all eddies subgrid &
\C{A1}{mynnrun}{a taper $\Psi_{bl}(2.5\dx/z_i)$ on $L$ and on the plume size $\le 1.2\dx$; equals 0.99 at 500~m for $z_i = 1$~km, so inert here} &
\C{A1}{ysu}{no $\Delta$ anywhere} &
\C{A1}{shh}{partition functions $P_{NL}, P_L(\Delta/z_i)$ scale the nonlocal flux and $K$ towards zero as $\Delta \to 0$; the scheme switches itself off in the LES limit} &
\C{A1}{boulac}{no $\Delta$ anywhere} \\
\rowshade{A2}\rid{A2} ensemble mean &
\C{A2}{mynnrun}{ensemble closure} & \C{A2}{ysu}{ensemble closure} &
\C{A2}{shh}{ensemble closure with LES-fitted partitions} & \C{A2}{boulac}{ensemble closure} \\
\rowshade{B1}\rid{B1} no horizontal shear production &
\C{B1}{mynnrun}{two vertical shear terms} & \C{B1}{ysu}{no production term at all: first order} &
\C{B1}{shh}{as YSU} & \C{B1}{boulac}{vertical shear only, $K[(\pd{u}{z})^2 + (\pd{v}{z})^2]$} \\
\rowshade{B2}\rid{B2} no $W$ gradients &
\C{B2}{mynnrun}{$W$ absent} & \C{B2}{ysu}{$W$ absent} & \C{B2}{shh}{$W$ absent} & \C{B2}{boulac}{$W$ absent} \\
\rowshade{B3}\rid{B3} no horizontal flux divergence &
\C{B3}{mynnrun}{vertical divergence only; horizontal mixing by the numerical operator} & \C{B3}{ysu}{as MYNN} & \C{B3}{shh}{as MYNN} & \C{B3}{boulac}{as MYNN} \\
\rowshade{B4}\rid{B4} TKE column-local &
\C{B4}{mynnrun}{$\qsq$ advection exists as an option and is off in the control} & \C{B4}{ysu}{no TKE} & \C{B4}{shh}{no TKE} &
\C{B4}{boulac}{$e$ not advected, by design; vertical diffusion only} \\
\rowshade{B5}\rid{B5} mixing separable &
\C{B5}{mynnrun}{$K_v$ from TKE, Smagorinsky $K_h$ from numerics; no energy exchange} & \C{B5}{ysu}{as MYNN} & \C{B5}{shh}{as MYNN} & \C{B5}{boulac}{as MYNN} \\
\rowshade{B6}\rid{B6} flat MOST surface &
\C{B6}{mynnrun}{MOST surface layer; wall value $\qsq \propto u_*^2$} &
\C{B6}{ysu}{MOST; optionally subgrid-orography drag on the surface stress (\texttt{topo\_wind})} &
\C{B6}{shh}{MOST} &
\C{B6}{boulac}{MOST in the mean equations only: the TKE equation has no surface source, the lowest level receives TKE by transport from above} \\
\rowshade{B7}\rid{B7} model surfaces horizontal &
\C{B7}{mynnrun}{no horizontal derivative} & \C{B7}{ysu}{none} & \C{B7}{shh}{none} & \C{B7}{boulac}{none} \\
\rowshade{C1}\rid{C1} local down-gradient &
\C{C1}{mynnrun}{K-form plus a mass-flux plume model for heat, moisture and momentum (nonlocal)} &
\C{C1}{ysu}{K-form plus a countergradient term $\gamma_\theta/h$ and an entrainment flux at the top $\propto (z/h)^3$} &
\C{C1}{shh}{K-form plus a prescribed nonlocal heat-flux profile weighted by $P_{NL}$} &
\C{C1}{boulac}{K-form plus the Troen--Mahrt countergradient $\gamma = 10\,\avg{w\theta}_0/(w_* h)$} \\
\rowshade{C2}\rid{C2} one $K$ &
\C{C2}{mynnrun}{one $K_M$, one $K_H$, vertical only} & \C{C2}{ysu}{one $K_m$, $K_h = K_m/\Pr$} & \C{C2}{shh}{as YSU} & \C{C2}{boulac}{one $K$ for everything, $\Pr = 1$} \\
\rowshade{C3}\rid{C3} flux anti-parallel &
\C{C3}{mynnrun}{yes} & \C{C3}{ysu}{yes} & \C{C3}{shh}{yes} & \C{C3}{boulac}{yes} \\
\rowshade{C4}\rid{C4} eddy size from the vertical structure &
\C{C4}{mynnrun}{NN09 blend of $L_S$, $L_T$, $L_B$ with bounds on $L_T$, a BouLac blend above the PBL and the LES-limit taper} &
\C{C4}{ysu}{no length scale: a prescribed $K$ profile shaped by $h$ from the bulk Richardson number} &
\C{C4}{shh}{as YSU} &
\C{C4}{boulac}{how far a parcel with the local TKE travels up and down against the stratification: $l_k = \min(l_{up}, l_{down})$; no ground distance, no $z_i$} \\
\rowshade{C5}\rid{C5} one scale for $K$ and $\varepsilon$ &
\C{C5}{mynnrun}{$\varepsilon = q^3/(B_1 L)$ with the same $L$} & \C{C5}{ysu}{no dissipation modelled} & \C{C5}{shh}{no dissipation modelled} &
\C{C5}{boulac}{two scales: $l_k = \min(l_{up}, l_{down})$ mixes, $l_\varepsilon = \sqrt{l_{up} l_{down}}$ dissipates} \\
\rowshade{C6}\rid{C6} vertical Ri; cutoff &
\C{C6}{mynnrun}{no critical Richardson number: $A_2 \to A_2/(1 + \Ri)$ keeps $S_M, S_H > 0$ at all $\Ri$, a Prandtl-number cap bounds the ratio} &
\C{C6}{ysu}{bulk Richardson number sets $h$; above it a local scheme with $K_h \propto (1 + 5\Ri)^{-2}$, no cutoff} &
\C{C6}{shh}{as YSU} &
\C{C6}{boulac}{no stability functions: stability enters only through the parcel length} \\
\rowshade{C7}\rid{C7} local equilibrium &
\C{C7}{mynnrun}{yes} & \C{C7}{ysu}{yes, first order: everything diagnostic} & \C{C7}{shh}{yes} & \C{C7}{boulac}{yes, 1.5 order} \\
\rowshade{C8}\rid{C8} buoyancy via $\avg{w\theta_v}$ &
\C{C8}{mynnrun}{$P_b$ through $\avg{w\theta_v}$, plus plume TKE production} & \C{C8}{ysu}{no TKE budget; buoyancy through $h$ and $w_*$} & \C{C8}{shh}{as YSU} &
\C{C8}{boulac}{$-\beta g K (\pd{\theta}{z} - \gamma)$, dry $\theta$} \\
\rowshade{C9}\rid{C9} universal constants &
\C{C9}{mynnrun}{NN09 set with Kitamura's $C_2, C_3$} & \C{C9}{ysu}{$b = 6.8$, $\Pr_0 = 1 + 0.272$, 15.9 (Noh)} &
\C{C9}{shh}{YSU constants plus LES-fitted partition coefficients} & \C{C9}{boulac}{$C_k = 0.4$, $C_\varepsilon = 1/1.4$, $\gamma$-coefficient 10} \\
\rowshade{C10}\rid{C10} pressure terms &
\C{C10}{mynnrun}{$K_q = 3K_m$ for the TKE transport; Rotta behind $S_M, S_H$} & \C{C10}{ysu}{none modelled} & \C{C10}{shh}{none modelled} &
\C{C10}{boulac}{TKE diffused with $K$ itself; no pressure--strain model (1.5 order)} \\
\end{xltabular}
\landscapeclose

% ---------------------------------------------------------------------------
%  Table I.2b  ideas, 3D family
% ---------------------------------------------------------------------------
\landscapeopen
\subsection*{Table I.2b --- How each scheme handles the assumption: the 3D family}
{\footnotesize The idea only; constants, equation numbers and code lines are in Part~II. \fork\ marks where the WRF port or this project's fork differs from the publication. \par}
\vspace{0.5ex}
\scriptsize
\setlength{\tabcolsep}{3.5pt}
\renewcommand{\arraystretch}{1.1}
\begin{xltabular}{\linewidth}{@{}>{\raggedright\arraybackslash}p{2.0cm}
  >{\columncolor{tGXVIII}\raggedright\arraybackslash}X
  >{\columncolor{tGXIX}\raggedright\arraybackslash}X
  >{\raggedright\arraybackslash}X
  >{\columncolor{tAPPROX}\raggedright\arraybackslash}X
  >{\columncolor{tFULL}\raggedright\arraybackslash}X
  >{\raggedright\arraybackslash}X >{\raggedright\arraybackslash}X@{}}
\toprule
 & \hd{gxviii} & \hd{gxix} & \hdn{Ito et al.\ 2015} & \hd{approx} & \hd{full} & \hdn{SMS-3DTKE} & \hdn{LES limit} \\
\midrule
\endfirsthead
\toprule
 & \hd{gxviii} & \hd{gxix} & \hdn{Ito et al.\ 2015} & \hd{approx} & \hd{full} & \hdn{SMS-3DTKE} & \hdn{LES limit} \\
\midrule
\endhead
\bottomrule
\endlastfoot
\rowshade{A1}\rid{A1} all eddies subgrid &
\C{A1}{gxviii}{$\dx$ enters as the horizontal eddy length, not as a partition: grid-aware, not scale-aware} &
\C{A1}{gxix}{no $\Delta$: the horizontal length is a flow scale} &
\C{A1}{ito}{every length scale is multiplied by $F(\Delta/h) = P_{TKE}^{3/2}$, because $\varepsilon$ is filter-independent while $q^3$ is not} &
\C{A1}{approx}{a Honnert taper on $l$ in this fork\fork\ (the publication refuses any $\Delta$ dependence by design); inert at 500~m} &
\C{A1}{full}{as 3D-APPROX} &
\C{A1}{sms}{partition functions of $\Delta_h/z_i$ blend the LES and the mesoscale limit of $K_v$, of $\varepsilon$ and of the nonlocal terms} &
\C{A1}{les}{the opposite limit: every energetic eddy is resolved and the grid is the eddy size} \\
\rowshade{A2}\rid{A2} ensemble mean &
\C{A2}{gxviii}{ensemble} & \C{A2}{gxix}{ensemble} &
\C{A2}{ito}{a spatial top-hat filter, declared as such} &
\C{A2}{approx}{ensemble, declared: this is why $l \neq f(\Delta)$} & \C{A2}{full}{as 3D-APPROX} &
\C{A2}{sms}{filter closure in the LES limit, ensemble in the mesoscale limit} & \C{A2}{les}{filter closure} \\
\rowshade{B1}\rid{B1} no horizontal shear production &
\C{B1}{gxviii}{the four horizontal gradients of the horizontal wind enter a Smagorinsky-form source $(0.2\dx)^2|D_h|^3$; scalars untouched} &
\C{B1}{gxix}{the same four gradients with the flow length $U^2 T_{L,u} T_{L,v}$ in place of $(0.2\dx)^2$} &
\C{B1}{ito}{all nine gradients in K-form stresses $\tau_{ij} = -L q S_M\,\pd{U_i}{x_j}$} &
\C{B1}{approx}{the $\qsq$ budget keeps the two vertical terms; the stresses are built under the boundary-layer approximation} &
\C{B1}{full}{all nine gradients in the stresses and in $P_s$: the algebraic $10 \times 10$ system} &
\C{B1}{sms}{all nine deformation components, $K_{mh}$ on the horizontal ones and $K_{mv}$ on those with $\pd{}{z}$} &
\C{B1}{les}{all nine, as SMS-3DTKE} \\
\rowshade{B2}\rid{B2} no $W$ gradients &
\C{B2}{gxviii}{no $W$ gradient in the source} & \C{B2}{gxix}{as G18} &
\C{B2}{ito}{$\pd{W}{x_j}$ in the K-form} &
\C{B2}{approx}{$\pd{W}{z}$ dropped by fiat} &
\C{B2}{full}{$\pd{W}{x}$, $\pd{W}{y}$, $\pd{W}{z}$ in all six stress rows} &
\C{B2}{sms}{$D_{13}, D_{23}, D_{33}$} & \C{B2}{les}{as SMS-3DTKE} \\
\rowshade{B3}\rid{B3} no horizontal flux divergence &
\C{B3}{gxviii}{unchanged: Smagorinsky numerics} & \C{B3}{gxix}{unchanged} &
\C{B3}{ito}{3D divergence of all K-form fluxes in the dynamical core} &
\C{B3}{approx}{3D divergence of all six stresses and of the heat and moisture fluxes, including a tendency on $W$} &
\C{B3}{full}{as 3D-APPROX} &
\C{B3}{sms}{$K_{mh}$, $K_{hh}$ horizontal diffusion, metric-corrected} & \C{B3}{les}{as SMS-3DTKE} \\
\rowshade{B4}\rid{B4} TKE column-local &
\C{B4}{gxviii}{$\qsq$ advected by WRF on the 500~m domain; no turbulent horizontal diffusion of its own} &
\C{B4}{gxix}{as G18} &
\C{B4}{ito}{3D transport with the horizontal length $L_x$} &
\C{B4}{approx}{$\qsq$ advected and filtered by WRF; the closure's own horizontal $\qsq$ diffusion runs only in 3D-FULL\fork} &
\C{B4}{full}{advected, plus horizontal diffusion $S_q l q$} &
\C{B4}{sms}{advected; diffused with $2K_{mh}$ horizontally} & \C{B4}{les}{as SMS-3DTKE} \\
\rowshade{B5}\rid{B5} mixing separable &
\C{B5}{gxviii}{the deformation now feeds $\qsq$, but momentum is still mixed by the numerical operator, with a different invariant: a one-way energy source} &
\C{B5}{gxix}{as G18} &
\C{B5}{ito}{one closure for both directions, with $L$ vertical and $L_x$ horizontal} &
\C{B5}{approx}{\texttt{diff\_opt=0}: the closure supplies all subgrid mixing} & \C{B5}{full}{as 3D-APPROX} &
\C{B5}{sms}{one prognostic $e$ feeds $K_{mh}$ and $K_{mv}$; the Smagorinsky part fades towards the LES limit} &
\C{B5}{les}{one $e$, one closure} \\
\rowshade{B6}\rid{B6} flat MOST surface &
\C{B6}{gxviii}{MOST surface layer} & \C{B6}{gxix}{MOST} & \C{B6}{ito}{bulk MOST} &
\C{B6}{approx}{MOST: $\avg{uw} = -u_*^2$ at the ground face; the 3D surface layer of Eghdami et al.\ exists as an option and is off\fork} &
\C{B6}{full}{as 3D-APPROX} & \C{B6}{sms}{MOST} & \C{B6}{les}{MOST as wall model} \\
\rowshade{B7}\rid{B7} model surfaces horizontal &
\C{B7}{gxviii}{uses WRF's metric-corrected deformations of the previous time step} & \C{B7}{gxix}{as G18} &
\C{B7}{ito}{flat domain} &
\C{B7}{approx}{metric-corrected gradients $\pd{}{x}|_z = \pd{}{x}|_\eta - z_x\,\pd{}{z}$; horizontal tendencies damped by a slope factor $\sigma$\fork} &
\C{B7}{full}{as 3D-APPROX; the horizontal production is paired with the damped tendency\fork} &
\C{B7}{sms}{\texttt{diff\_opt=2} metric correction; an ad hoc slope taper on the Smagorinsky part} &
\C{B7}{les}{metric correction, no taper} \\
\rowshade{C1}\rid{C1} local down-gradient &
\C{C1}{gxviii}{K-form} & \C{C1}{gxix}{K-form} & \C{C1}{ito}{K-form plus the level-3 countergradient $\Gamma_\theta$} &
\C{C1}{approx}{K-form for the vertical fluxes; the other twelve moments are algebraic by-products} &
\C{C1}{full}{no $K$: the stresses solve an algebraic system coupling all gradients} &
\C{C1}{sms}{K-form plus a prescribed nonlocal heat-flux profile and a countergradient momentum term, vertical only} &
\C{C1}{les}{K-form} \\
\rowshade{C2}\rid{C2} one $K$ &
\C{C2}{gxviii}{horizontal $K_h \neq$ vertical $K_v$: $l_h = 0.2\dx$} &
\C{C2}{gxix}{two horizontal scales $T_{L,u}, T_{L,v}$ and a vertical one} &
\C{C2}{ito}{$L_x \neq L$, and $S_M = S_H = 1$ horizontally} &
\C{C2}{approx}{one $l$; the anisotropy of the variances is diagnosed as a perturbation of $\qsq/3$} &
\C{C2}{full}{the tensor is solved, its principal axes are free; one master $l$ remains} &
\C{C2}{sms}{$K_{mh}(\Delta_h) \neq K_{mv}(l_v)$; horizontal Prandtl number 1/3} &
\C{C2}{les}{$K_{mh} \neq K_{mv}$ on an anisotropic grid} \\
\rowshade{C3}\rid{C3} flux anti-parallel &
\C{C3}{gxviii}{yes} & \C{C3}{gxix}{yes} & \C{C3}{ito}{yes} &
\C{C3}{approx}{vertical fluxes aligned; $\avg{uv}$, $\avg{u\theta}$ are products of vertical fluxes and vertical gradients} &
\C{C3}{full}{no: counter-gradient stress components and locally negative production are admissible} &
\C{C3}{sms}{yes} & \C{C3}{les}{yes} \\
\rowshade{C4}\rid{C4} eddy size from the vertical structure &
\C{C4}{gxviii}{vertical: NN09; horizontal: $0.2\dx$} &
\C{C4}{gxix}{horizontal: $U\,T_L$ with $T_L = 0.15\, z_i/\sigma_{u,v}$; $z_i$ from a rewritten boundary-layer-height routine\fork, variances from flat-terrain profiles} &
\C{C4}{ito}{NN09 $L \times F(\Delta/h)$; horizontal $L_x = 0.1\,h\,F$} &
\C{C4}{approx}{Blackadar $l = \kappa z l_0/(\kappa z + l_0)$ with $l_0$ from the whole column; Deardorff limit $0.53\,q/N$ and Honnert taper\fork} &
\C{C4}{full}{as 3D-APPROX, plus the strain cap $l \le 12\,q/(B_1|S|)$\fork} &
\C{C4}{sms}{$l_h = \sqrt{\dx\Delta y}$; $l_v$ blends $\min(\Delta z, 0.76\sqrt{e}/N)$ with an NN09-type mesoscale length} &
\C{C4}{les}{$l = \min(\Delta,\ 0.76\sqrt{e}/N)$} \\
\rowshade{C5}\rid{C5} one scale for $K$ and $\varepsilon$ &
\C{C5}{gxviii}{same $L$ in $\varepsilon$; the horizontal source is dissipated with the vertical scale} & \C{C5}{gxix}{as G18} &
\C{C5}{ito}{the same $F$ on every length; a locality option makes $L_\varepsilon \propto q$} &
\C{C5}{approx}{same $l$ (the BouLac option would separate them)} &
\C{C5}{full}{same $l$, including the strain-capped one\fork} &
\C{C5}{sms}{a separately blended dissipation length; $C_\varepsilon$ grows with $l/\Delta$} &
\C{C5}{les}{$C_\varepsilon = c_{\varepsilon1} + c_{\varepsilon2}\,l/\Delta$} \\
\rowshade{C6}\rid{C6} vertical Ri; cutoff &
\C{C6}{gxviii}{a source independent of Ri, added to a scheme that keeps its cutoff} & \C{C6}{gxix}{as G18} &
\C{C6}{ito}{NN09 $S_M, S_H$ unchanged} &
\C{C6}{approx}{level-2 cutoff $\Rfc = 0.191$ (MY82 constants) in $S_M, S_H$} &
\C{C6}{full}{no $S_M, S_H$: the equilibrium follows from the full strain; $N$ enters through $l$ only} &
\C{C6}{sms}{no stability functions; $N$ limits $l$; no cutoff} & \C{C6}{les}{as SMS-3DTKE} \\
\rowshade{C7}\rid{C7} local equilibrium &
\C{C7}{gxviii}{yes} & \C{C7}{gxix}{yes} & \C{C7}{ito}{yes} & \C{C7}{approx}{yes} &
\C{C7}{full}{yes --- the strain cap is the admission that it fails under fast strain\fork} &
\C{C7}{sms}{yes, 1.5 order: $K$ is instantaneous} & \C{C7}{les}{yes} \\
\rowshade{C8}\rid{C8} buoyancy via $\avg{w\theta_v}$ &
\C{C8}{gxviii}{$P_b$ through $\avg{w\theta_v}$} & \C{C8}{gxix}{as G18} &
\C{C8}{ito}{horizontal heat fluxes $-L_x q\,\pd{\theta}{x}$ exist} &
\C{C8}{approx}{buoyancy redistributes between $\avg{w^2}$ and the horizontal variances through $\avg{w\theta_v}$; $\avg{u\theta_v}$ diagnosed from vertical gradients} &
\C{C8}{full}{$\beta g_j \avg{u_i\theta_v}$ in every stress row and three heat-flux components; the $\qsq$ budget still uses $\avg{w\theta_v}$ only} &
\C{C8}{sms}{horizontal heat flux $K_{hh}\,\pd{\theta}{x}$ exists; $P_b$ vertical} & \C{C8}{les}{as SMS-3DTKE} \\
\rowshade{C9}\rid{C9} universal constants &
\C{C9}{gxviii}{NN09 set plus $c = 0.2$} & \C{C9}{gxix}{NN09 set plus 0.15 (Hanna)} &
\C{C9}{ito}{NN09 set unchanged} &
\C{C9}{approx}{MY82 set in this fork, BouLac-tuned set in the publication\fork} & \C{C9}{full}{as 3D-APPROX} &
\C{C9}{sms}{$c_k = 0.15$, $c_s = 0.25$; the code's mesoscale constants differ from the paper's} &
\C{C9}{les}{$c_k = 0.15$, $c_s = 0.25$} \\
\rowshade{C10}\rid{C10} pressure terms &
\C{C10}{gxviii}{as MYNN} & \C{C10}{gxix}{as MYNN} & \C{C10}{ito}{Rotta with $F L$} &
\C{C10}{approx}{$S_q l q$ transport; Rotta plus the $C_1$ isotropisation} & \C{C10}{full}{as 3D-APPROX} &
\C{C10}{sms}{$2K\,\pd{e}{x_i}$ in three dimensions} & \C{C10}{les}{as SMS-3DTKE} \\
\end{xltabular}
\landscapeclose
